\documentclass[10pt,a4paper]{amsart}
\usepackage[margin=19mm]{geometry}
\usepackage{amsmath,amssymb,mathtools}
\usepackage[scr=rsfs,cal=euler]{mathalfa}
\usepackage{xfrac}
\usepackage[unicode,hidelinks,bookmarks=false]{hyperref}
\newcommand{\ie}{i.e.\ }
\newcommand{\cf}{cf.\ }
\newcommand{\resp}{respectively\ }
\newcommand{\Subsection}{\subsection}
\providecommand{\nobreakdash}{\nobreak\hbox{-}}
\setlength{\emergencystretch}{2em}
\multlinegap0pt
\usepackage{amssymb}
\newtheorem{theorem}{Theorem}
\title{Maximal subalgebras of the Lie algebra $W_n(\mathbb{K})$}
\author{Y.Chapovskyi, A.Petravchuk, O.Tyshchenko}
\date{Source: arXiv:2605.22970v1 (CC BY 4.0)}
\begin{document}
\maketitle

\noindent\emph{Source and licence.} Maximal subalgebras of the Lie algebra $W_n(\mathbb{K})$. Version arXiv:2605.22970v1 (CC BY 4.0), distributed by arXiv under the Creative Commons Attribution 4.0 International licence, http://creativecommons.org/licenses/by/4.0/, as recorded in the arXiv licence metadata for this exact version and shown on the abstract page https://arxiv.org/abs/2605.22970v1. Full licence notice: \texttt{licenses/arxiv-2605.22970-CC-BY-4.0.txt}.\par\smallskip

\noindent\textit{Editorial scope.} Counted unit: the article's theorem theorem (source0.tex lines 485--490). The article's complete original proof of it is delivered with it (source0.tex lines 492--523, one \texttt{proof} environment of 32 lines). No mathematics is written, repaired or re-derived: the statement and the proof are the article's own lines, in the article's own notation, and no retained line is substituted or re-flowed.  No further source-local context is needed: every cross-reference of every retained line resolves inside the two delivered spans.  Nothing was omitted from any retained span, so the package carries no omission record.  No retained line contains a citation, so the wrapper carries no bibliography and no bibliography entry.  No retained line contains a figure environment, an image-inclusion command or drawing code, so no image is delivered and the package declares no asset dependency.  Editorial formatting only, in addition: an AMS \texttt{amsart} preamble that restates the article's own declarations and macros used by the retained lines, namely the environments \texttt{theorem} on separate counters, together with the standard packages those lines need.  The retained environments print in this document as Theorem 1 in order of appearance, whereas the article prints the same items with the numbering of its own full document.  The complete native source of the article as distributed in the arXiv e-print for this exact version is bundled unchanged as \texttt{sources/201173/source0.tex}.
	\begin{theorem}
		Let $I_k = (x_{k+1}, \dots, x_n)$ be the ideal in $A$ generated by the last $n-k$ variables. Denote by $L_k$ the subalgebra that preserves the ideal $I_k$, i.e., 
		$L_k = \{ D \in W_n \mid D(I_k) \subseteq I_k \}.$
		Then $L_k$ is a maximal subalgebra of $W_n$ for $i=1, \ldots , n-1.$ 
		
	\end{theorem}

	\begin{proof} One can easily to show that 
		$$L_k=A\frac{\partial}{\partial x_1} + \dots +A\frac{\partial}{\partial x_k} + I_k \frac{\partial}{\partial x_{k+1}} + \dots + I_k \frac{\partial}{\partial x_n}.
		$$
		Since $L_k$ is a Lie subalgebra in $W_n$ we can consider $L_k$  and $W_n$ as   $L_k$-modules.  
		Consider the quotient module $Q=W_n/L_k.$ Denote ${\frac{\overline{\partial}}{\partial x_i}}=\frac{\partial}{\partial x_i}+L_k, i=1, \ldots , n.$ One can easily  show that 
		$$			Q = W_n / L_k=(A/I_k){\frac{\overline{\partial}}{\partial x_{k+1}}} + \dots + (A/I_k){\frac{\overline{\partial}}{\partial x_n}}.$$
		Note that the quotient ring $A/I_k $ is isomorphic to the polynomial ring $K[x_1, \dots, x_k],$ so we have  
		$$W_n/L_k\simeq {K}[x_1, \dots, x_k]{\frac{\overline{\partial}}{\partial x_{k+1}}}+\cdots + {K}[x_1, \dots, x_k]{\frac{\overline{\partial}}{\partial x_n}}.
		$$
		Our aim is to show that $Q=W_n/L_k$ is an irreducible $L_k$-module. 
		Take any $\overline{D}\in Q.$	Then $\overline{D}$ is of the form  
		$$		
		\overline{D} = f_{k+1}(x_1, \dots, x_k){\frac{\overline{\partial}}{\partial x_{k+1}}} + \cdots + f_n(x_1, \dots, x_k){\frac{\overline{\partial}}{\partial x_n}}.
		$$	
		Since  $\frac{\partial}{\partial x_i} \in L_k$ for $1 \le i \le k$, we have:
		$$[\frac{\partial}{\partial x_i},  \overline{D}] = \frac{\partial f_{k+1}}{\partial x_i} \frac{\overline{\partial}}{\partial x_{k+1}} + \dots + \frac{\partial f_n}{\partial x_i} \frac{\overline{\partial}}{\partial x_n} \in L_k\langle \overline{D}\rangle .
		$$
		$$	\overline{D}_1 = c_{k+1}\frac{\overline{\partial}}{\partial x_{k+1}} + \dots + c_n \frac{\overline{\partial}}{\partial x_n} \in L_k\langle \overline{D}\rangle ,
		$$
		where $c_{k+1}, \dots, c_n \in K$ and $c_j \ne 0$ for some $j,$ $k+1\leq j\leq n.$
		
		Then $[-x_j \frac{\partial}{\partial x_j}, \overline{D}_1] = c_j \frac{\overline{\partial}}{\partial x_j} \in L_k \langle \overline{D} \rangle$.
		Finally, 
		\begin{equation*}
			[-x_j f(x_1, \dots, x_k) \frac{\partial}{\partial x_q}, \frac{\overline{\partial}}{\partial x_j}] = f(x_1, \dots , x_k) \frac{\partial}{\partial x_q} \in L_k \langle \overline{D} \rangle,
		\end{equation*} 
		for all $q = k+1,\dots,n$, $f \in \mathbb{K}[x_1, \dots, x_k]$.
		
		Thus, $L_k \langle \overline{D} \rangle = Q$ for an arbitrary nonzero $\overline{D} \in Q, $ which means $Q$ is an irreducible $L_k$-module.
		It follows directly that $L_k$ is a maximal subalgebra in $W_n$. 
		Indeed, if there existed an intermediate subalgebra $L_k \subsetneq M \subsetneq W_n$, then there would exist a non-trivial $L_k$-submodule $M/L_k$ in $W_n/L_k = Q$, which contradicts the irreducibility of the module $Q$.
	\end{proof}

\end{document}
